\subsection{Collective Coherent Phases of the \texorpdfstring{$U(1)$}{U(1)} Fiber}
  \label{subsec:collective-U1-coherence}

  The previous sections described an effective $U(1)$ fiber sector of the non-injective projection
  of the $\chi$ substrate, with an integer fiber winding whose identification with electric charge
  is proposed but not derived (Section~\ref{sec:charge-quantisation-pi1}).
  We now consider the possibility that this fiber sector admits collective coherent phases.

  The analysis below does not introduce new fundamental fields.
  It studies a specific collective regime of already established degrees of freedom.
  Superconductivity is interpreted as a macroscopic coherent phase of the $U(1)$ fiber,
  arising from the stabilization and phase locking of composite topological excitations.

  \subsubsection*{Stability of \texorpdfstring{$w=2$}{w=2} Composite Classes}
    \label{subsubsec:w2-stability}

    Section~\ref{sec:particles-as-localized-excitations-of-the-chi-field} reads charged excitations as
    topological configurations of the $\chi$ field projected onto the $U(1)$ fiber.
    The association of a single excitation of winding number $w=1$ with the effective charge $e$ is a
    proposed physical identification, $q_{\mathrm{eff}} = w \cdot e$, which is not derived
    (Section~\ref{sec:charge-quantisation-pi1}).

    We now consider composite configurations of winding number $w=2$.
    Such configurations are not trivial superpositions of two independent $w=1$ modes.
    They define a distinct topological class characterized by a continuous fiber matching structure.

    Let $L_{\Pi}$ denote the effective projective ordering operator restricted to the fiber sector.
    We define the projective stability gap
    \begin{equation}
      \Delta_{\Pi} \equiv \lambda_{\min}\!\left(L_{\Pi}\big|_{w=2}\right),
    \end{equation}
    as the lowest positive eigenvalue around the $w=2$ configuration.

    A finite $\Delta_{\Pi}$ implies that the composite class is metastable against
    separation into two $w=1$ excitations or local phase defects.
    A positive gap is an additional assumption of this collective model.
    Section~\ref{sec:necessity-noninjective} does not derive it from bounded transport capacity;
    a specified operator, spectrum and stability analysis are still required.

  \subsubsection*{Phase Coherence and the London Structure}
    \label{subsubsec:london-structure}

    If a macroscopic density of $w=2$ composites occupies a common fiber phase sector,
    we may introduce the order parameter
    \begin{equation}
      \Psi = |\Psi| e^{i\theta},
    \end{equation}
    where $|\Psi|^2$ measures the density of coherent composites.

    Two hypotheses are needed here and are not derived: the proposed identification
    $q_{\mathrm{eff}} = w \cdot e$, so that a $w=2$ composite carries charge $2e$, and a $U(1)$
    potential~$A$ under which $\Psi$ transforms with that charge, $\theta \to \theta + 2e\lambda/\hbar$ when
    $A \to A + \nabla\lambda$.
    Under them, gauge covariance fixes the unique first-order covariant gradient
    \begin{equation}
      D\theta = \nabla\theta - \frac{2e}{\hbar} A.
    \end{equation}

    The leading contribution to the free energy must then take the form
    \begin{equation}
      F_{\text{grad}} =
      \frac{\hbar^2 \rho_s}{2 m^*}
      \left(\nabla\theta - \frac{2e}{\hbar} A\right)^2,
    \end{equation}
    where $m^*$ is an effective inertia and $\rho_s$, which has the dimension of a number density in this
    normalisation, sets the phase stiffness $\hbar^2\rho_s/m^*$.

    This term is not postulated as a symmetry-breaking potential.
    It represents the energetic cost of breaking global fiber coherence.
    Under the same hypotheses, with the supercurrent defined as
    \begin{equation}
      j_s = -\frac{\partial F_{\text{grad}}}{\partial A}
        = \frac{2e\hbar\rho_s}{m^*}\left(\nabla\theta - \frac{2e}{\hbar} A\right),
    \end{equation}
    this relation between the supercurrent and~$A$ is the London relation.
    Flux quantization follows on a closed path~$C$ under three further conditions: $\rho_s > 0$ at every point
    of~$C$, the phase~$\theta$ well defined along~$C$ (that is, $\Psi \neq 0$ there and $\Psi$ single-valued), and
    $j_s = 0$ on~$C$.
    Then $\nabla\theta = (2e/\hbar)A$ on~$C$ and $\oint_C \nabla\theta \cdot d\ell = 2\pi n$ with $n \in \mathbb{Z}$,
    so the magnetic flux through a surface bounded by~$C$ is $n\,h/2e$.
    Where $\rho_s = 0$, $j_s = 0$ holds trivially and nothing is quantized.

  \subsubsection*{Frustration-Induced Stabilization in Strongly Correlated Regimes}
    \label{subsubsec:frustration-stabilization}

    In conventional low-$T_c$ materials, mediator-assisted interactions lower
    $\Delta_{\Pi}$ and stabilize the $w=2$ class.

    In strongly correlated systems, stabilization may arise instead from
    real-space frustration in admissibility constraints of the fiber.
    Let $F$ denote a frustration density proxy, operationally linked to
    staggered magnetic correlations.

    In this regime, isolated $w=1$ excitations increase local mismatch,
    whereas a $w=2$ composite redistributes phase incompatibility
    over an extended region.
    Pair formation thus minimizes a geometric cost functional rather than
    resulting from a specific momentum-space interaction kernel.

    The amplitude scale for composite formation is set by the dominant
    frustration energy scale $J$,
    \begin{equation}
      \Delta_{\Pi}^{\text{ampl}} \sim J.
    \end{equation}

  \subsubsection*{Phase Stiffness and BKT-Type Transition}
    \label{subsubsec:phase-stiffness}

    Global superconductivity requires phase locking across the system.
    The relevant control parameter is the phase stiffness $\rho_s$,
    which measures the admissible projective capacity of the projected substrate.
    In the two relations of this paragraph, $\rho_s$ denotes the two-dimensional phase stiffness in energy
    units; in the normalisation of $F_{\text{grad}}$, where $\rho_s$ is a density, it is
    $\hbar^2\rho_s^{(2\mathrm{D})}/m^*$ with $\rho_s^{(2\mathrm{D})}$ the areal density.

    In doped quasi-two-dimensional systems, the density of active composites
    scales with carrier concentration $\delta$.
    To leading order,
    \begin{equation}
      \rho_s(0) \propto \delta J.
    \end{equation}

    In strongly anisotropic layered systems, the transition temperature
    is governed by vortex unbinding and satisfies approximately
    \begin{equation}
      k_B T_c \simeq \frac{\pi}{2} \rho_s(T_c^-).
    \end{equation}

    This naturally allows separation between amplitude formation
    ($T^* \sim \Delta_{\Pi}^{\text{ampl}}/k_B$)
    and global phase coherence ($T_c$).

  \subsubsection*{Superconductivity as Collective Saturation}
    \label{subsubsec:sc-collective-saturation}

    The emergence of superconductivity can thus be interpreted as a collective
    saturation phase of the $U(1)$ fiber.

    At the microscopic level, bounded projective transport constrains
    admissible field gradients.
    At the macroscopic level, coherent occupation of a single fiber phase
    sector is read as minimizing projective ordering cost; this reading is
    interpretive.
    Flux expulsion requires more than the hypotheses of
    Section~\ref{subsec:collective-U1-coherence}.
    In a region where the regime is static, $\rho_s > 0$ is uniform and no vortex is present (so that
    $\nabla \times \nabla\theta = 0$), and with $A$ obeying the static Maxwell equation
    $\nabla \times B = \mu_0 j_s$ sourced by the supercurrent of
    Section~\ref{subsec:collective-U1-coherence}, the curl of the London relation gives
    $\nabla \times j_s = -(4e^2\rho_s/m^*)\,B$, hence
    \begin{equation}
      \nabla^2 B = \frac{B}{\lambda_L^2},
      \qquad
      \lambda_L^2 = \frac{m^*}{4\mu_0 e^2 \rho_s},
    \end{equation}
    with $\rho_s$ normalised as in $F_{\text{grad}}$.
    Magnetic flux is then expelled from the bulk of that region over the penetration depth~$\lambda_L$ (Meissner
    effect).
    Neither the Maxwell dynamics of~$A$ nor the coupling hypotheses are derived here.

    In this perspective, superconductivity is structurally analogous to
    other saturation phenomena derived in this work.
    Mass corresponds to isotropic saturation of local projective capacity.
    Schwinger pair production corresponds to saturation of directed admissible flux.
    Gravitational curvature corresponds to collective constraint on the
    effective projective ordering rate.

    Superconductivity represents the coherent phase of the same bounded
    projective dynamics in the electromagnetic sector.

    No modification of Standard Model interactions is required.
    The phenomenon emerges as a specific collective regime of the
    already established $U(1)$ fiber structure.
